% TP2, Autómatas celulares off-lattice. Estilo tomado de la presentación del TP5.
% Compilar: pdflatex document.tex (dos veces). Figuras: venv/bin/python presentacion/generar_figuras.py
% Fuente de datos: tp2_resultados/ (sweep_summary.csv, datos_caracteristicos/, benchmark/). No se re-simuló.
% Los MP4 se reproducen sobre el póster en un visor compatible (Okular, Acrobat).
% Para entregar: \entregatrue y completar las URLs públicas (YouTube/Vimeo).
\documentclass[aspectratio=169,20pt]{beamer}
\geometry{paperwidth=320mm,paperheight=180mm}
\usepackage[utf8]{inputenc}
\usepackage[T1]{fontenc}
\usepackage{lmodern}
\usefonttheme{professionalfonts}
\usefonttheme[onlymath]{serif}
\usepackage[spanish,es-nodecimaldot]{babel}
\usepackage{amsmath,amssymb,graphicx,booktabs,tikz,multimedia}
\usetikzlibrary{arrows.meta,positioning,calc}
\usetheme{Frankfurt}
\usecolortheme{beaver}
\definecolor{darkred}{RGB}{153,0,0}
\definecolor{VI}{RGB}{31,119,180}
\definecolor{VO}{RGB}{214,39,40}
\setbeamercolor{frametitle}{bg=gray!8,fg=darkred}
\setbeamercolor{title}{fg=darkred}
\setbeamercolor{structure}{fg=darkred}
\setbeamercolor{section in head/foot}{bg=darkred,fg=white}
\setbeamertemplate{navigation symbols}{}
\setbeamertemplate{blocks}[default]
\setbeamersize{text margin left=12mm,text margin right=12mm}
\setbeamerfont{frametitle}{size=\fontsize{28}{33}\selectfont}
\setbeamerfont{title}{size=\fontsize{34}{40}\selectfont}
\newif\ifentrega\entregafalse
% Links públicos de las animaciones (completar para la entrega)
\newcommand{\urlVI}{}
\newcommand{\urlVO}{}
\setbeamertemplate{headline}{\begin{beamercolorbox}[wd=\paperwidth,ht=7mm,dp=2mm,leftskip=12mm]{section in head/foot}\fontsize{15}{18}\selectfont\insertsectionhead\end{beamercolorbox}}
\setbeamertemplate{footline}{\hbox{\begin{beamercolorbox}[wd=.85\paperwidth,ht=6mm,dp=2mm,leftskip=12mm]{section in head/foot}\fontsize{13}{16}\selectfont Grupo 11\quad Autómatas celulares off-lattice\end{beamercolorbox}\begin{beamercolorbox}[wd=.15\paperwidth,ht=6mm,dp=2mm,rightskip=8mm plus1fil]{section in head/foot}\fontsize{13}{16}\selectfont\insertpagenumber/\inserttotalframenumber\end{beamercolorbox}}}
\setbeamertemplate{itemize items}[circle]
\setlength{\parskip}{4mm}
\newcommand{\vb}[1]{\mathbf{#1}}
\newcommand{\notaBreve}[1]{{\fontsize{18}{22}\selectfont #1}}
% Figura (2/3 del ancho) con condiciones fijas al costado: {título}{archivo}{texto lateral}
\newcommand{\grafico}[3]{\begin{frame}{#1}\begin{columns}[c]
 \begin{column}{.68\textwidth}\centering\includegraphics[width=\textwidth,height=.80\textheight,keepaspectratio]{figuras/#2}\end{column}
 \begin{column}{.30\textwidth}\notaBreve{#3}\end{column}\end{columns}\end{frame}}
% Animación: póster en el PDF de entrega con link; MP4 embebido en la versión para exponer.
\newcommand{\anim}[3]{\ifentrega\includegraphics[height=.62\textheight]{figuras/#1.png}\par
 \ifx#3\empty\notaBreve{\textcolor{magenta}{(insertar link)}}\else\notaBreve{\expandafter\url\expandafter{#3}}\fi
 \else\movie[poster,showcontrols]{\includegraphics[height=.62\textheight]{figuras/#1.png}}{figuras/#1.mp4}\fi\par\notaBreve{#2}}
% Comparación Vicsek | Votante con los mismos ejes: {título}{prefijo}{parámetros}
\newcommand{\comparacion}[3]{\begin{frame}{#1}\centering
 \begin{tabular}{@{}l@{\hspace{8mm}}l@{}}\hspace{20mm}\textcolor{VI}{Vicsek}&\hspace{14mm}\textcolor{VO}{Votante}\\
 \includegraphics[height=.68\textheight]{figuras/#2_vicsek.png}&\includegraphics[height=.68\textheight]{figuras/#2_voter.png}\end{tabular}
 \par\vspace{1mm}\notaBreve{#3}\end{frame}}
% Dos animaciones lado a lado: {título}{a}{rótulo a}{url a}{b}{rótulo b}{url b}{parámetros}
\newcommand{\animpar}[8]{\begin{frame}{#1}\begin{columns}[T]
 \begin{column}{.49\textwidth}\centering\anim{#2}{#3}{#4}\end{column}
 \begin{column}{.49\textwidth}\centering\anim{#5}{#6}{#7}\end{column}\end{columns}
 \vfill\notaBreve{#8}\end{frame}}
\tikzset{bloque/.style={draw=darkred,thick,align=center,minimum height=15mm,text width=72mm},flecha/.style={->,very thick,>=Stealth}}
\AtBeginSection[]{\begin{frame}\centering\usebeamerfont{title}\textcolor{darkred}{\insertsectionhead}\end{frame}}
\title{Autómatas celulares off-lattice}
\subtitle{Bandadas de agentes autopropulsados: Vicsek y votante}
\author{Grupo 11\\Santiago Diaz -- 63058\quad Lautaro Gazzaneo -- 61484\\Luca Maximiliano Rossi -- 63730}
\institute{Simulación de Sistemas -- ITBA}
\date{18 de noviembre de 2026}
\begin{document}
\begin{frame}\titlepage\end{frame}

\section{Introducción}
\begin{frame}{Sistema real: bandadas}
\begin{columns}[c]
\begin{column}{.48\textwidth}
\begin{itemize}\setlength{\itemsep}{3mm}
\item Aves, cardúmenes, bacterias, robots
\item Cada agente se mueve a rapidez constante
\item Solo interactúa con vecinos cercanos
\item Emerge (o no) un movimiento colectivo ordenado
\end{itemize}
\end{column}
\begin{column}{.48\textwidth}
\[
 \vb x_i(t+\Delta t)=\vb x_i(t)+\vb v_i(t+\Delta t)\,\Delta t
\]
\[
 \vb v_i=v_0\,(\cos\theta_i,\ \sin\theta_i)
\]
\[
 \mathcal N_i(t)=\{j\neq i:\ \|\vb x_i-\vb x_j\|\le r_c\}
\]
\end{column}
\end{columns}
\end{frame}

\begin{frame}{Reglas de actualización del ángulo}
\begin{columns}[T]
\begin{column}{.49\textwidth}
\textcolor{VI}{Vicsek: promedio de vecinos}
\[
 \theta_i(t+\Delta t)=\operatorname{atan2}\!\Big(\langle\sin\theta\rangle_{r_c},\,\langle\cos\theta\rangle_{r_c}\Big)+\Delta\theta_i
\]
\notaBreve{$\langle\cdot\rangle_{r_c}$: promedio sobre $\mathcal N_i\cup\{i\}$.}
\end{column}
\begin{column}{.49\textwidth}
\textcolor{VO}{Votante: copia de un vecino}
\[
 \theta_i(t+\Delta t)=\theta_{j}(t)+\Delta\theta_i,\qquad j\sim\mathcal U(\mathcal N_i)
\]
\notaBreve{Si $\mathcal N_i=\varnothing$, conserva su ángulo: $\theta_j=\theta_i$.}
\end{column}
\end{columns}
\vspace{8mm}
\[
 \Delta\theta_i\sim\mathcal U\!\left[-\tfrac{\eta}{2},\,\tfrac{\eta}{2}\right]\qquad\eta:\ \text{amplitud del ruido}
\]
\end{frame}

\section{Implementación}
\begin{frame}{Arquitectura del motor}
\centering
\begin{tikzpicture}
\node[bloque] (sys) at (0,0) {CellularAutomataSystem\\\notaBreve{step(), observables}};
\node[bloque] (rule) at (16,0) {UpdateRule\\\notaBreve{computeNewAngle()}};
\node[bloque] (cim) at (0,-4) {ParticleSystem\\\notaBreve{Cell Index Method}};
\node[bloque] (p) at (0,-8) {Particle\\\notaBreve{$\vb x$, $\theta$, $\vb v$, vecinos}};
\node[bloque,text width=48mm] (vi) at (12.5,-4.5) {\textcolor{VI}{VicsekRule}};
\node[bloque,text width=48mm] (vo) at (19.5,-4.5) {\textcolor{VO}{VoterRule}};
\draw[flecha] (sys)--(rule);
\draw[flecha] (sys)--(cim);
\draw[flecha] (cim)--(p);
\draw[flecha,dashed] (vi.north)--(vi.north |- rule.south);
\draw[flecha,dashed] (vo.north)--(vo.north |- rule.south);
\end{tikzpicture}
\end{frame}

\begin{frame}{Un paso de simulación}
\begin{columns}[c]
\begin{column}{.55\textwidth}
\begin{enumerate}\setlength{\itemsep}{3mm}
\item CIM con contorno periódico: $\mathcal N_i$ para todo $i$
\item Para cada $i$: $\theta_i^{\,\mathrm{nuevo}}=\text{regla}(\theta_i,\mathcal N_i,\eta)$
\item Recién entonces, para cada $i$:
\[ \theta_i\leftarrow\theta_i^{\,\mathrm{nuevo}},\quad \vb x_i\leftarrow(\vb x_i+\vb v_i\Delta t)\bmod L \]
\item Registrar $v_a(t)$ y $S(t)$
\end{enumerate}
\end{column}
\begin{column}{.42\textwidth}\centering
\begin{tikzpicture}[x=9mm,y=9mm]
\foreach \i in {0,...,5}{\draw[gray!60] (\i,0)--(\i,5);\draw[gray!60] (0,\i)--(5,\i);}
\fill[darkred!25] (2,2) rectangle (3,3);
\fill[darkred!10] (2,3) rectangle (4,4); \fill[darkred!10] (3,2) rectangle (4,3); \fill[darkred!10] (3,1) rectangle (4,2);
\draw[thick] (0,0) rectangle (5,5);
\fill[darkred] (2.5,2.5) circle (.08);
\draw[dashed,darkred,thick] (2.5,2.5) circle (1);
\draw[<->] (0,-.4)--node[below]{$L/M\ge r_c$}(1,-.4);
\end{tikzpicture}
\par\notaBreve{Celda propia y 4 vecinas; pares guardados en ambos sentidos.}
\end{column}
\end{columns}
\end{frame}

\section{Simulaciones}
\begin{frame}{Sistema y parámetros}
\begin{columns}[c]
\begin{column}{.40\textwidth}\centering
\begin{tikzpicture}[x=11mm,y=11mm]
\draw[very thick,dashed] (0,0) rectangle (6,6);
\foreach \x/\y/\a in {1/4/20,2/2/40,4/5/-10,5/3/30,3/4/10,4.5/1.5/60,1.5/1/0,3.2/2.8/25}{\draw[-{Stealth},VI,thick] (\x,\y)--++(\a:.6);}
\draw[darkred,thick] (3.2,2.8) circle (1);
\draw[darkred] (3.2,2.8)--node[above,pos=.7]{$r_c$}++(200:1);
\draw[<->] (-.5,0)--node[left]{$L$}(-.5,6);
\node[below] at (3,-.1){\notaBreve{Contorno periódico}};
\end{tikzpicture}
\end{column}
\begin{column}{.57\textwidth}
{\renewcommand{\arraystretch}{1.6}
\begin{tabular}{@{}ll@{}}
\multicolumn{2}{@{}l}{\textcolor{darkred}{Fijos}}\\
$L=10$ & $v_0=0{,}03$\\
$r_c=1$ & $\Delta t=1$\\
2000 pasos & \\[4mm]
\multicolumn{2}{@{}l}{\textcolor{darkred}{Variables}}\\
Modelo & Vicsek, votante\\
$\rho$ & $\frac{1}{3\pi},\ \frac{1}{2\pi},\ \frac{1}{\pi},\ 2,\ 4,\ 8$\\
$N$ & $10,\ 15,\ 31,\ 200,\ 400,\ 800$\\
$\eta$ & $0;\ 0{,}5;\ \ldots;\ 5$ rad\\
\end{tabular}}
\end{column}
\end{columns}
\end{frame}

\begin{frame}{Observables}
\[
 v_a(t)=\frac{1}{N v_0}\left\|\sum_{i=1}^{N}\vb v_i(t)\right\|\qquad
 S(t)=\frac{N_{\max}(t)}{N}
\]
\notaBreve{$N_{\max}$: partículas del mayor cluster (unidas por cadenas de vecinos a distancia $\le r_c$)}
\[
 \overline{v_a}^{(k)}=\frac{1}{T-t_e}\sum_{t=t_e}^{T-1}v_a^{(k)}(t)\qquad
 \langle v_a\rangle=\frac1M\sum_{k=1}^{M}\overline{v_a}^{(k)}\qquad
 \sigma=\sqrt{\frac{1}{M-1}\sum_{k=1}^{M}\left(\overline{v_a}^{(k)}-\langle v_a\rangle\right)^2}
\]
\notaBreve{$T=2000$}
\end{frame}

\section{Resultados}
\subsection{Polarización}
\animpar{Animaciones}{anim_vicsek_etabaja}{\textcolor{VI}{Vicsek}}{\urlVI}{anim_voter_etabaja}{\textcolor{VO}{Votante}}{\urlVO}{$\rho=4$\qquad$\eta=0{,}5$ rad}

\comparacion{Evolución del parámetro de orden}{evol_va}{$\eta=1$ rad\qquad Línea vertical: inicio del estacionario}

\comparacion{Parámetro de orden vs. ruido}{va_vs_eta}{$M=5$ realizaciones}


\subsection{Clusters}

\comparacion{Evolución de la fracción del cluster mayor: $\rho=2,4,8$}{evol_s_altas}{$\eta=1$ rad\qquad Línea vertical: inicio del estacionario}

\comparacion{Evolución de la fracción del cluster mayor: $\rho<1$}{evol_s_bajas}{$\eta=1$ rad\qquad Línea vertical: inicio del estacionario}

\comparacion{Fracción del cluster mayor vs. ruido: $\rho=2,4,8$}{s_vs_eta_altas}{$M=5$ realizaciones}

\comparacion{Fracción del cluster mayor vs. ruido: $\rho<1$}{s_vs_eta_bajas}{$M=5$ realizaciones}

\subsection{Polarización y clusters}
\comparacion{Parámetro de orden vs. fracción del cluster mayor: $\rho=2,4,8$}{va_vs_s_altas}{$M=5$ realizaciones}

\comparacion{Parámetro de orden vs. fracción del cluster mayor: $\rho<1$}{va_vs_s_bajas}{$M=5$ realizaciones}

\subsection{Tiempos del CIM}
\grafico{Cell Index Method: tiempo de búsqueda}{benchmark.png}{$L=20$\par $13\times13$ celdas\par 5 repeticiones}

\section{Conclusiones}
\begin{frame}{Conclusiones}
\begin{itemize}\setlength{\itemsep}{4mm}
\item Vicsek mantiene el orden hasta ruidos mucho mayores que el votante
\item Al aumentar la densidad, Vicsek se ordena más y el votante menos
\item A baja densidad, el ruido fragmenta el grupo y el orden cae junto con el cluster mayor
\item A alta densidad, el sistema queda siempre conectado: se desordena sin fragmentarse
\item El CIM es un orden de magnitud más rápido que fuerza bruta y sus tiempos son similares a los del TP1
\end{itemize}
\end{frame}

\begin{frame}\centering\usebeamerfont{title}\textcolor{darkred}{Muchas gracias}\end{frame}
\end{document}
